\chapter{Regurgitation}

\section*{Definition}
Regurgitation is the effortless return of ingested food or liquid from the esophagus or stomach into the mouth or pharynx, distinct from vomiting which involves active contraction of abdominal muscles. It may occur immediately after eating or be delayed by hours.

\section*{What Happens in Your Body}
Regurgitation results from failure of the anti-reflux barrier at the gastroesophageal junction. The lower esophageal sphincter (LES) typically maintains a high-pressure zone; transient LES relaxations, hypotensive LES, or anatomic disruption (hiatal hernia) permit retrograde flow. In rumination syndrome, habitual contraction of abdominal muscles actively returns gastric contents to the mouth without nausea.

\section*{Causes}
\begin{itemize}
  \item Gastroesophageal reflux disease (GERD -- most common cause)
  \item Hiatal hernia (sliding or paraesophageal)
  \item Rumination syndrome (behavioral, often involuntary)
  \item Achalasia (failure of LES relaxation with esophageal dilatation)
  \item Esophageal stricture or web (mechanical obstruction)
  \item Esophageal diverticulum (Zenker diverticulum -- food trapped in pouch)
  \item Gastroparesis (delayed gastric emptying -- diabetes, post-surgical)
  \item Pregnancy (hormonal LES relaxation and increased intra-abdominal pressure)
  \item Obesity (increased intra-abdominal pressure)
  \item Medications: Nitrates, calcium channel blockers, anticholinergics (LES relaxation)
\end{itemize}

\section*{When to See a Doctor}
See your doctor if:
\begin{itemize}
  \item Regurgitation with difficulty swallowing (dysphagia) or painful swallowing (odynophagia)
  \item Regurgitation of undigested food hours after eating (suggests diverticulum or achalasia)
  \item Unexplained weight loss or anemia
  \item Regurgitation with coughing, choking, or recurrent pneumonia (aspiration)
  \item Regurgitation in infants with failure to thrive
  \item Hematemesis (blood in regurgitated material -- consider esophageal tear)
\end{itemize}

\section*{Related Symptoms}
\begin{itemize}
  \item Heartburn
  \item Dysphagia
  \item Chest pain
  \item Cough (especially nocturnal)
  \item Hoarseness
  \item Sore throat
  \item Dental erosion
  \item Aspiration
  \item Halitosis
  \item Belching
\end{itemize}

\section*{Self-Care / Home Management}
\begin{itemize}
  \item Eat smaller, more frequent meals and avoid lying down within 2-3 hours after eating.
  \item Elevate head of bed 6-8 inches with blocks or a wedge pillow.
  \item Avoid trigger foods: spicy, fatty, acidic foods, chocolate, caffeine, alcohol.
  \item Weight loss if overweight.
  \item Over-the-counter antacids, H2 receptor antagonists, or PPIs for symptom relief.
  \item Avoid tight-fitting clothing around the abdomen.
\end{itemize}

\section*{How Your Doctor Will Evaluate You}
\begin{itemize}
  \item Upper endoscopy (EGD) to evaluate esophageal and gastric mucosa.
  \item Esophageal pH monitoring to confirm GERD and quantify acid exposure.
  \item Esophageal manometry to assess LES pressure and peristalsis.
  \item Barium swallow study to identify structural abnormalities.
  \item Gastric emptying study if gastroparesis suspected.
\end{itemize}

\section*{Treatment Options}
\begin{itemize}
  \item Lifestyle modifications and pharmacotherapy (PPIs are first-line for GERD).
  \item Antireflux surgery (Nissen fundoplication) for refractory GERD.
  \item Botulinum toxin injection into LES for achalasia.
  \item Pneumatic dilation or Heller myotomy for achalasia.
  \item Cognitive behavioral therapy or diaphragmatic breathing for rumination syndrome.
  \item Prokinetic agents for gastroparesis (metoclopramide, domperidone).
\end{itemize}

\section*{References}
\begin{thebibliography}{99}
\bibitem{ref1} Katz PO, Dunbar KB, Schnoll-Sussman FH, et al. "ACG Clinical Guideline for the Diagnosis and Management of Gastroesophageal Reflux Disease." \textit{American Journal of Gastroenterology}, 2022;117(1):27--56.
\bibitem{ref2} Kahrilas PJ, Pandolfino JE. "Treatments for gastroesophageal reflux disease." \textit{JAMA}, 2018;319(9):924--925.
\bibitem{ref3} Vaezi MF, Pandolfino JE, Vela MF. "ACG clinical guideline: Diagnosis and management of achalasia." \textit{American Journal of Gastroenterology}, 2013;108(8):1238--1249.
\bibitem{ref4} Tack J, Talley NJ, Camilleri M, et al. "Functional gastroduodenal disorders." \textit{Gastroenterology}, 2006;130(5):1466--1479.
\bibitem{ref5} Murray HB, Juarascio AS, Di Lorenzo C, et al. "Diagnosis and treatment of rumination syndrome: A critical review." \textit{American Journal of Gastroenterology}, 2019;114(4):562--578.
\bibitem{ref6} Richter JE, Rubenstein JH. "Presentation and epidemiology of gastroesophageal reflux disease." \textit{Gastroenterology}, 2018;154(2):267--276.
\end{thebibliography}
